\mysection{Duffのデバイス}
\myindex{Duff's device}
\myindex{Unrolled loop}

Duffのデバイス
\footnote{\href{http://go.yurichev.com/17137}{wikipedia}}
は、その中央にジャンプすることが可能なアンロールドループです。
アンロールドループは、fallthrough switch()文を使って実装されています。
ここでは、Tom Duffの元のコードを少し簡素化したバージョンを使用します。
メモリ内の領域をクリアする関数を書かなければならないとしましょう。
1バイトずつクリアする単純なループを考えることができます。
これは明らかに遅いです。なぜなら、すべての現代のコンピュータはより広いメモリバスを持っているからです。
そのため、4バイトまたは8バイトのブロックを使用してメモリ領域をクリアする方が良い方法です。
ここでは64ビットの例を扱うため、8バイトブロックでメモリをクリアすることにします。
ここまでは順調です。
しかし、末尾はどうでしょうか？
メモリクリアルーチンは、8の倍数でないサイズの領域に対しても呼び出される可能性があります。
そこで、アルゴリズムは次のようになります：

\begin{itemize}
\item 8バイトブロックの数を計算し、8バイト（64ビット）メモリアクセスを使用してそれらをクリアする

\item 末尾のサイズを計算し、1バイトメモリアクセスを使用してそれをクリアする
\end{itemize}

2番目のステップは、単純なループを使って実装できます。
しかし、アンロールドループとして実装してみましょう：

\lstinputlisting[style=customc]{\CURPATH/duff_EN.c}

まず、計算がどのように実行されるかを理解しましょう。
メモリ領域のサイズは64ビット値として渡されます。
そして、この値は2つの部分に分割できます：

% .... 7 6 5 4 3 2 1 0
%|....|B|B|B|B|B|S|S|S|

\begin{center}
\begin{bytefield}[endianness=big,bitwidth=0.03\linewidth]{11}
\bitheader{7,6,5,4,3,2,1,0} \\
\bitbox{3}{\dots} & 
\bitbox{1}{B} & 
\bitbox{1}{B} & 
\bitbox{1}{B} & 
\bitbox{1}{B} & 
\bitbox{1}{B} & 
\bitbox{1}{S} & 
\bitbox{1}{S} & 
\bitbox{1}{S}
\end{bytefield}
\end{center}

（\q{B}は8バイトブロック数、\q{S}は末尾のバイト単位での長さ）。

入力メモリ領域サイズを8で割ると、値は単に3ビット右にシフトされます。
しかし、余りを計算するには、最下位3ビットを分離するだけで済みます！
そのため、8バイトブロック数は$count>>3$として計算され、余りは$count \& 7$として計算されます。
また、8バイト処理を実際に実行するかどうかを調べる必要があるため、$count$の値が7より大きいかどうかをチェックする必要があります。
これは、最下位3ビットをクリアし、結果の数をゼロと比較することで行います。なぜなら、ここで必要なのは「$count$の上位部分がゼロでないか」という質問に答えることだからです。
もちろん、これは8が$2^{3}$であり、$2^n$の数による除算が簡単だからこそ動作します。
他の数では不可能です。
実際に、これらのハックを使う価値があるかどうかは言い難いです。なぜなら、読みにくいコードになるからです。
しかし、これらのトリックは非常に人気があり、実践的なプログラマーは、たとえ使用しなくても、それらを理解する必要があります。

そのため、最初の部分は単純です：8バイトブロックの数を取得し、64ビットゼロ値をメモリに書き込みます。
2番目の部分は、fallthrough switch()文として実装されたアンロールドループです。

まず、ここで行わなければならないことを平易な英語で表現しましょう。

「$count\&7$の値が教えてくれる分だけ、メモリにゼロバイトを書き込む」必要があります。
それが0なら、終了にジャンプします。やる仕事はありません。
それが1なら、switch()文内の場所にジャンプし、そこで1つのストア操作のみが実行されます。
それが2なら、別の場所にジャンプし、そこで2つのストア操作が実行されます。以下同様です。
入力値7は、7つすべての操作の実行につながります。
8はありません。なぜなら、8バイトのメモリ領域は関数の最初の部分によって処理されるからです。
そのため、アンロールドループを書きました。
これは古いコンピュータでは確実に通常のループより速かったです（逆に、最新の\ac{CPU}は短いループの方がアンロールされたものより良く動作します）。
おそらく、これは現代の低コスト組込み\ac{MCU}でまだ意味があるかもしれません。

最適化するMSVC 2012が何をするか見てみましょう：

\lstinputlisting[style=customasmx86]{\CURPATH/duff_MSVC2012_x64_Ox_EN.asm}

関数の最初の部分は予測可能です。
2番目の部分は、単にアンロールドループと、その中の正しい命令に制御フローを渡すジャンプです。
\TT{MOV}/\TT{INC}命令ペアの間には他のコードがないため、実行は必要な数のペアを実行して最後まで続きます。
ちなみに、\TT{MOV}/\TT{INC}ペアが固定バイト数（3+3）を消費することがわかります。
そのため、ペアは6バイトを消費します。
これを知っていれば、switch()ジャンプテーブルを除去できます。入力値に6を乗じて、$current\_RIP + input\_value * 6$にジャンプするだけです。

これも高速化できる可能性があります。なぜなら、ジャンプテーブルから値を取得する必要がないからです。

6はおそらく高速乗算にとって非常に良い定数ではなく、価値がないかもしれませんが、アイデアは理解できるでしょう\footnote{練習として、ジャンプテーブルを除去するためにコードを書き直してみることができます。
\myindex{x86!\Instructions!STOSB}
命令ペアは4バイトまたは8バイトを消費するように書き直すことができます。
1バイトも可能です（\TT{STOSB}命令を使用）}。

これは、昔のデモ制作者がアンロールドループで行っていたことです。

\subsection{アンロールドループを使うべきか？}
\myindex{Unrolled loop}

アンロールドループは、\ac{RAM}と\ac{CPU}の間に高速キャッシュメモリがなく、\ac{CPU}が次の命令のコードを取得するために毎回\ac{RAM}からロードしなければならない場合に利点があります。
これは現代の低コスト\ac{MCU}と古い\ac{CPU}の場合です。

アンロールドループは、\ac{RAM}と\ac{CPU}の間に高速キャッシュがあり、ループの本体がキャッシュに収まる場合は短いループより遅くなります。\ac{CPU}は\ac{RAM}に触れることなくそこからコードをロードします。
高速ループとは、本体のサイズがL1キャッシュに収まるループですが、さらに高速なループは、マイクロオペレーションキャッシュに収まる小さなループです。